\section{The actual differential exceptional locus of the point chart}

Throughout this chapter $\pi$ is the existing actual polynomial point
chart map on the original $\mathbb C^n$,
\[
 \pi_e(z)=z_e,\qquad \pi_i(z)=z_e(b_i+z_i)\quad(i\ne e).
\]
The exceptional locus means the original locus where the actual
holomorphic manifold differential is not bijective. No properness or
global blowup universal property is assumed for this single chart map.

\begin{theorem}[The actual manifold differential of the genuine point chart]
\label{hir:actual-point-blowup-differential-coordinates}
\uses{hir:actual-analytic-blowup-coordinate-pullback-taylor-naturality}
\lean{Hironaka.HironakaActualPointBlowupExceptionalLocus.mfderiv_actual_point_blowup_apply}
\leanok
For every point $z$, vector $v$ and coordinate $i$, the actual manifold
holomorphic differential of the unchanged point chart satisfies
\[
 (D\pi_z(v))_i=
 \begin{cases}
 v_e,&i=e,\\
 v_e(b_i+z_i)+z_e v_i,&i\ne e.
 \end{cases}
\]
\end{theorem}
\begin{proof}
The original point chart is genuinely holomorphic everywhere. Compose
its actual Fr\'echet differential with the continuous complex-linear
$i$-th coordinate projection. The chain rule identifies this composition
with the derivative of its actual scalar coordinate function. For
$i=e$ that function is the original linear coordinate projection. For
$i\ne e$ the product rule applied to $z_e(b_i+z_i)$ gives the displayed
expression. Uniqueness of the Fr\'echet differential proves equality for
every vector. On the original normed-space manifold with its self model,
the actual manifold differential is the actual Fr\'echet differential;
this existing dictionary transfers the computed expression to the
original manifold differential, rather than postulating a Jacobian.
\end{proof}

\begin{theorem}[The genuine coordinate hyperplane lies in the actual exceptional locus in dimension at least two]
\label{hir:actual-point-blowup-hyperplane-exceptional}
\uses{hir:actual-point-blowup-differential-coordinates}
\lean{Hironaka.HironakaActualPointBlowupExceptionalLocus.exists_actual_point_blowup_nonzero_kernel,Hironaka.HironakaActualPointBlowupExceptionalLocus.mem_exceptionalLocus_actual_point_blowup_of_coordinate_eq_zero}
\leanok
Suppose $n\ge2$ and $z_e=0$. The actual manifold differential of $\pi$
at $z$ has a nonzero kernel vector. Consequently $z$ belongs to the
existing differential-defined exceptional locus of $\pi$.
\end{theorem}
\begin{proof}
Choose a genuine coordinate $j\ne e$, possible because $n\ge2$, and let
$v$ be the actual standard coordinate vector $e_j$. Its $j$-th coordinate
is one, so $v\ne0$, whereas $v_e=0$. The preceding differential formula
gives zero in output coordinate $e$. Every other output coordinate is
$v_e(b_i+z_i)+z_e v_i=0$. Thus the whole actual differential sends the
nonzero vector $v$ to zero. It is not injective and therefore not
bijective. This is exactly membership in the original exceptional locus.
\end{proof}

\begin{theorem}[The actual one-dimensional point chart is the identity]
\label{hir:actual-point-blowup-one-dimensional-identity}
\uses{hir:actual-analytic-blowup-coordinate-pullback-taylor-naturality}
\lean{Hironaka.HironakaActualPointBlowupExceptionalLocus.actual_point_blowup_one_eq_id,Hironaka.HironakaActualPointBlowupExceptionalLocus.exceptionalLocus_actual_point_blowup_one,Hironaka.HironakaActualPointBlowupExceptionalLocus.strictTransform_actual_point_blowup_one}
\leanok
In dimension one, for every actual chart coordinate and parameter, the
unchanged point chart is the identity. Its original exceptional locus
is empty. For every closed set $A$, its original geometric strict
transform under this chart map is exactly $A$.
\end{theorem}
\begin{proof}
There is only one coordinate, namely $e$, so each component of $\pi(z)$
is the corresponding component of $z$. Hence the actual map is the
identity. Its actual manifold differential is the identity linear map,
which is bijective at every point; its original exceptional locus is
empty. The existing strict-transform definition is the closure of the
inverse image of $A$ outside that exceptional locus. It consequently
reduces to $\overline A=A$ because $A$ is closed.
\end{proof}

These statements supply the differential-defined exceptional-set input
for the later point geometric comparison. They do not equate the
normalized-member ideal with an original geometric strict-transform
ideal. In dimension one that proposed equality can fail, so the genuine
identity case above is retained separately. General positive-dimensional
permissible centres require their own compatibility argument.
